% paper_v12.tex  — major revision from v11
% v12 changes (vs v11):
%   - No page limit
%   - Full-width figures (figure* or large figures) with bigger fonts
%   - New §5.4: SSL failure section — domain-specific SSL v1+v2, coding ceiling
%   - SSL + verifier unified story: feature-space failure, not algorithm failure
%   - Section 5 restructured: "Five Methods, One Ceiling" framing
%   - Four new figures replacing the single small fig_knot_domain_profiles_v3
%   - Abstract rewritten around convergent-evidence framing
%   - Title updated to reflect multi-method scope
\documentclass[10pt,twocolumn]{article}
\usepackage[margin=0.75in,top=0.85in,bottom=0.85in]{geometry}
\usepackage[T1]{fontenc}
\usepackage[utf8]{inputenc}
\usepackage{microtype}
\usepackage{amsmath,amssymb}
\usepackage{booktabs,tabularx,array,multirow}
\usepackage{graphicx}
\usepackage{xcolor}
\usepackage[numbers,sort&compress]{natbib}
\usepackage{hyperref}
\hypersetup{colorlinks=true,linkcolor=blue!50!black,citecolor=blue!50!black,urlcolor=blue!50!black}
\setlength{\parindent}{0.9em}
\setlength{\parskip}{0.2em}
\renewcommand{\arraystretch}{1.1}
\newcolumntype{C}{>{\centering\arraybackslash}p{1.6cm}}
\graphicspath{{../figures/}{./}}

\title{\textbf{Code Correctness Is Not in the Text:\\[4pt]
\large Five Methods, One Ceiling — Convergent Evidence for\\[2pt]
CoT Quality Measurement Breakdown in Coding}}
\author{Yuhan Chi\\Fudan University\\\texttt{masterwuguicyh@gmail.com}}
\date{}

\begin{document}
\maketitle

% ─────────────────────────────────────────────────────────────────────────────
\begin{abstract}
% ─────────────────────────────────────────────────────────────────────────────

Can the quality of a coding solution be read from its chain-of-thought (CoT) trace? We test this question using five methods on the same feature family—token-level uncertainty signals and text-trajectory statistics—across three domains: math competitions, science (GPQA), and coding (LiveCodeBench-v5). All experiments use a single model, \textbf{DeepSeek-R1-0528-Qwen3-8B}.

In math and science, the answer is clearly \emph{yes}: a simple interpretable probe achieves AUC-of-AUROC 0.958 and 0.799 respectively. The signal is temporally structured and transfers across early-stopping anchor positions. A best-of-$N=64$ reranking experiment delivers $+10.0\%$ and $+8.0\%$ pass@1 gains.

In coding, the answer is consistently \emph{no}. (1)~The same probe yields AUC-of-AUROC 0.434, \emph{below} a single-feature token-confidence baseline. (2)~A sweep of 83 coding-specific scalars peaks at AUROC 0.556. (3)~A nonlinear MLP classifier does not escape the ceiling. (4)~Self-supervised pre-training on 42{,}000 unlabeled CoT traces—including a carefully engineered domain-specific variant with VICReg anti-collapse and five-term structured loss—tops out near 0.50 AUROC for coding at full labels. (5)~A token-level surface-noise ablation (de-knotting) with all per-token signal arrays properly masked shows that in math, knot tokens \emph{carry} the discriminative signal ($\Delta$AUROC $= -0.049$ after removal), while in coding, no masking strategy recovers signal.

These five methods converge on a single explanation: code correctness depends on executable state evolution, and this information is not encoded in token trajectories. No amount of data, no change of algorithm, and no surface cleaning recovers a signal that is absent from the feature family. The failure is a \emph{measurement} failure, not a modeling failure.

\end{abstract}

% ─────────────────────────────────────────────────────────────────────────────
\section{Introduction}
% ─────────────────────────────────────────────────────────────────────────────

Process supervision and verifier-based reranking rest on a shared assumption: that chain-of-thought traces encode answer quality, and that this quality signal can be extracted by a probe of reasonable complexity. This is well-established for mathematical reasoning \citep{cobbe2021training,lightman2023verify,wang2024mathshepherd}. The open question is whether the same is true across domains.

This paper is not about building a better verifier. It is about \emph{measuring} whether a specific feature family encodes correctness in each domain—and about accumulating enough evidence to be confident about a null result in coding.

Our starting point is deliberately simple: five hand-specified features (three trajectory statistics, one token-confidence summary, and one neuron-consistency mean), a \texttt{StandardScaler $\to$ LogisticRegression} pipeline, evaluated at anchor positions $\{10, 40, 70, 100\}$\% of the visible trace. Simple enough that a performance failure can be attributed to the features, not to the model. We then expand: 83 coding-side scalars, a nonlinear MLP, self-supervised pre-training with up to 42,000 unlabeled runs, and a token-level ablation study.

The central finding is a \textbf{convergent measurement breakdown}. Every method applied to the token-trajectory feature family reaches the same ceiling near AUROC 0.5 for coding correctness. In contrast, math and science show strong, temporally structured, transferable signals from the same features. The gap is not reducible to label scarcity, model capacity, or surface noise. The most parsimonious account is that code correctness depends on executable state evolution—an object that text-surface features observe only as an imperfect shadow.

A downstream best-of-$N=64$ reranking experiment confirms the practical consequence: the probe gains $+10.0\%$ pass@1 on math and $+8.0\%$ on science, but $-0.6\%$ (effectively zero) on coding.

% ─────────────────────────────────────────────────────────────────────────────
\section{Related Work}
% ─────────────────────────────────────────────────────────────────────────────

\paragraph{Process supervision and verifiers.} Learning outcome reward models \citep{cobbe2021training} and process reward models \citep{lightman2023verify} established that step-level supervision improves over outcome-level alone. Math-Shepherd \citep{wang2024mathshepherd} automated this with LLM annotation; ProcessBench \citep{zheng2025processbench} benchmarks step-level error detection; generative verifier formulations \citep{zhang2025genrm} frame reward modeling as next-token prediction. Our work does not compete with any of these—it asks whether simpler, cheaper features measure the same thing across domains.

\paragraph{CoT faithfulness and self-correction.} Whether CoT traces faithfully reflect model reasoning is contested \citep{lanham2023faithfulness}. Intrinsic self-correction work \citep{huang2024selfcorrect,liu2024intrinsic} documents that reflective text does not always improve reasoning; self-refine \citep{madaan2023selfrefine} and reflexion \citep{shinn2023reflexion} show that iterative revision helps only with external feedback. We add a cross-domain dimension: not only does self-correction fail to reliably improve reasoning, but the \emph{measurability} of reasoning quality from text trajectory is itself domain-conditioned.

\paragraph{Self-supervised and semi-supervised probing.} Self-supervised learning for representation quality has been studied in NLP \citep{devlin2019bert}, but its application to CoT trajectory probing is unexplored. We evaluate SSL pre-training directly as a label-efficiency tool—if the features encode something latent about correctness, SSL should help at low label fractions. The finding that SSL tops out near AUROC 0.5 for coding at full labels rules out ``need more labeled data'' as the bottleneck.

\paragraph{Measurement invariance.} The psychometric literature on measurement invariance \citep{vandenberg2000review} establishes that the same instrument can measure different constructs in different populations. Probing work in NLP \citep{belinkov2022probing} makes a related point. We extend this logic to verifier settings: a feature family that measures correctness in math may observe a qualitatively different object in coding.

% ─────────────────────────────────────────────────────────────────────────────
\section{Setup}
\label{sec:setup}
% ─────────────────────────────────────────────────────────────────────────────

\subsection{Model, domains, and data}

All CoT traces are generated by \textbf{DeepSeek-R1-0528-Qwen3-8B} at temperature 1.0.
\textbf{Math:} 7,680 aligned runs across AIME24, AIME25, BRUMO25, and HMMT25 (competition math; ground-truth binary correctness).
\textbf{Science:} 12,672 runs from GPQA (graduate-level science MCQ).
\textbf{Coding:} 10,688 runs from LiveCodeBench-v5 (competitive programming; judge-verified correctness).

\subsection{Feature family}

We use five hand-specified features computed from cached token streams and neuron activations. These are kept intentionally cheap so that failure is attributable to the feature family, not to engineering.

\textit{Token-level signals.} \texttt{tok\_conf} is the mean negative log probability over visible prefix tokens (per-token token confidence / entropy). Four additional arrays are cached per token: \texttt{tok\_gini}, \texttt{tok\_logprob}, \texttt{tok\_neg\_entropy}, \texttt{tok\_selfcert}. All five are available per-token and used in ablation studies.

\textit{Trajectory statistics.} \texttt{traj\_continuity}: mean Jaccard similarity between consecutive 32-token activation slices. \texttt{traj\_novelty}: mean $1 - \max_j \text{Jaccard}(s_i, s_j)$ over all prior slices. \texttt{traj\_reflection\_count}: negated count of non-adjacent slice revisits with Jaccard $> 0.30$.

\textit{Neuron summary.} \texttt{nc\_mean}: mean row width of the sparse activation bank across the visible prefix.

\subsection{Probe pipeline and metrics}

The base probe is \texttt{StandardScaler} $\to$ \texttt{TruncatedSVD} $\to$ \texttt{LogisticRegression}, fit separately per domain. Evaluation uses leave-one-problem-out cross-validation (GroupKFold by \texttt{problem\_id}).

The primary metric is \textbf{AUC-of-AUROC (AoA)}: area under the anchor-wise AUROC curve at positions $\{10, 40, 70, 100\}$\%, used as an early-stopping summary. Bootstrap 95\% CIs on AoA use $B=10{,}000$ resamplings of datasets within each domain (7 for math, 2 for science).

% ─────────────────────────────────────────────────────────────────────────────
\section{Math and Science: Strong, Temporally Structured, Transferable}
\label{sec:positive}
% ─────────────────────────────────────────────────────────────────────────────

Figure~\ref{fig:auroc} and Table~\ref{tab:positive} summarize the main probe result. Math achieves AoA 0.958 [0.931, 0.980] and AUROC 0.982 at the full trace. Science reaches AoA 0.799 [0.775, 0.822] and AUROC 0.841. Coding falls to AoA 0.434 and AUROC 0.407---below the token-confidence baseline (0.506).

\begin{figure*}[t]
  \centering
  \includegraphics[width=\textwidth]{fig_auroc_by_anchor_v12.pdf}
  \caption{%
    \textbf{(a)} AUROC at each anchor position (10\%–100\% of trace visible) for the three domains.
    Math and science rise with more context; coding oscillates near 0.5 regardless of how much
    trace is revealed. \textbf{(b)} AUC-of-AUROC (AoA) with 95\% bootstrap CIs. The domain gap
    is large and statistically well-separated.
  }
  \label{fig:auroc}
\end{figure*}

\begin{table}[h]
\centering
\small
\begin{tabular}{lcccc}
\toprule
Domain & AoA (95\% CI) & A@100 & Fwd $\Delta$ & Bwd $\Delta$ \\
\midrule
Math    & \textbf{.958} [.931, .980] & \textbf{.982} & $-$.001 & $-$.006 \\
Science & \textbf{.799} [.775, .822] & \textbf{.841} & $-$.010 & $-$.001 \\
Coding  &  .434 [.404, .464] & .407 & --- & --- \\
\midrule
\multicolumn{5}{l}{\textit{Baselines}} \\
\ \ \texttt{tok\_conf} & .506/.809/.506 & .490/.856/.490 & --- & --- \\
\ \ SC oracle$^\dagger$ & .953/.921/.961 & .953/.921/.961 & --- & --- \\
\bottomrule
\end{tabular}
\caption{Main holdout result. AoA = AUC-of-AUROC; 95\% CI from per-dataset bootstrap ($B=10{,}000$). A@100 = AUROC at 100\% anchor. Frozen-basis $\Delta$AUROC: Fwd = source-to-target transfer gap; Bwd = target-to-source. $^\dagger$SC oracle = self-consistency upper bound (requires ground truth).}
\label{tab:positive}
\end{table}

\paragraph{Temporal structure in math.} The math signal is not a final-score story—it is a \emph{process} story. The Spearman correlation between \texttt{traj\_reflection\_count} and correctness deepens as more of the trace is revealed:
\[
\rho_{10\%} = -0.306,\quad \rho_{40\%} = -0.383,\quad \rho_{100\%} = -0.491.
\]
Early reflection counts encode something real, but late reflection counts encode failure more strongly. This is the signature of a process feature: it measures whether the model converged, not just what it said at the end.

\paragraph{Frozen-basis transfer.} We freeze the low-rank basis learned at one anchor and refit only the logistic head at another: $\Delta\text{AUROC} = \text{AUROC}_{\text{frozen}} - \text{AUROC}_{\text{task-specific}}$. For math, forward and backward gaps are both near $-0.001$ to $-0.006$, indicating that the feature geometry is stable across anchor positions. Science has a larger forward gap ($-0.010$), suggesting reasoning maturity is less uniformly structured. Coding has no meaningful forward transfer to report because the within-domain AUROC is already near chance.

\paragraph{MLP sanity check.} Replacing the linear head with a one-hidden-layer MLP under five-fold GroupKFold gives $0.940/0.737/0.499$ mean-anchor AUROC for math/science/coding, and $0.965/0.760/0.507$ at the 100\% anchor. The cross-domain pattern is not a linear-head artifact.

% ─────────────────────────────────────────────────────────────────────────────
\section{Coding: Five Methods, One Ceiling}
\label{sec:coding}
% ─────────────────────────────────────────────────────────────────────────────

We apply five independent methods to ask whether anything in the token-trajectory feature family encodes coding correctness. All five reach the same answer.

\subsection{Method 1: Full-scale interpretable probe}

On the 10,688-run LiveCodeBench holdout table, the best pairwise proxy for self-corrective control-flow turbulence---\texttt{condition\_reversal\_per1k}---is statistically nonzero ($\rho = -0.062$, $p = 1.15\times 10^{-10}$). But the core probe achieves AoA 0.434, below the token-confidence baseline (0.506). The feature family carries the faintest trace of signal---but not enough to separate correct from incorrect runs usefully.

\subsection{Method 2: Broader feature sweep (83 scalars)}

We screen 83 interpretable coding-side scalars: struct complexity counts, control-flow features, comment density, CoT-trace length, code-trace alignment scores, and token-level statistics. The strongest negative scalar is \texttt{struct\_density} ($\rho = -0.110$, AUROC 0.565), followed by \texttt{cot\_trace\_count} ($\rho = -0.099$, AUROC 0.558). The strongest program-structure count \texttt{cf\_n\_ifs} has $\rho = +0.109$ (AUROC 0.563): more complex code is weakly \emph{more} correct, not less. Combining the top four interpretable scalars reaches only AUROC 0.556 under GroupKFold. No scalar or combination approaches the math probe level.

\subsection{Method 3: Nonlinear classifier (MLP)}

A one-hidden-layer MLP on the full 25-dimensional feature cache gives AUROC 0.499/0.507 at mean-anchor and 100\%-anchor respectively. Changing the classifier does not rescue coding. The ceiling is not a linear-model artifact.

\subsection{Method 4: Self-supervised pre-training}

\label{sec:ssl}
\paragraph{Motivation.} Methods 1--3 use labeled data exclusively. Method 4 asks a different question: if we train an encoder on all 42,000 \emph{unlabeled} CoT runs, can we learn a latent space where coding correctness becomes separable? If the answer were yes, it would imply that correctness information \emph{is} encoded in the feature space but is hard to extract with limited labeled data. If the answer is no, it implies the feature space simply does not contain the information.

\paragraph{SSL v1: shared cross-domain basis.} We train a masked-reconstruction SSL encoder on all 42,000 runs with ranks $r \in \{4, 8, 16\}$, using three loss terms: masked feature recovery, cross-anchor alignment, and raw/rank view consistency. Results:

\begin{itemize}
  \setlength\itemsep{0.15em}
  \item $r=8$: universal collapse to AUROC $= 0.5000$ on all 18 evaluation rows. The combined loss saturated at $\approx 0.432$ before epoch 60 and did not decrease through epoch 300. All samples collapsed to the same point in $\mathbb{R}^8$.
  \item $r=16$: partial recovery (AUROC 0.779 on math, 0.502 on coding at 100\% labels). Still well below the \texttt{no\_svd\_lr} baseline.
  \item No rank improved label efficiency in coding. At 1\% labels, \texttt{no\_svd\_lr} already outperforms or matches the best SSL variant at most fractions.
\end{itemize}

Failure modes: cross-domain pollution (math structure diluted by coding), reconstruction task too easy for linear autoencoder, no LR schedule to escape early saddle.

\paragraph{SSL v2: domain-specific structured SSL.} We redesign the SSL pipeline to address v1's failures: (1)~separate encoder per domain, (2)~five-term structured loss (family-block masking + NT-Xent contrastive + future-anchor prediction + raw/rank consistency + VICReg anti-collapse), (3)~cosine warm-restart LR schedule, (4)~up to 500 epochs. All 27 trained bundles (3 domains $\times$ 3 ranks $\times$ 3 seeds) achieve zero dead dimensions. Loss decreases by multiple orders of magnitude (math $r=8$: $325\text{k} \to 5\text{k}$; coding $r=24$: $208\text{k} \to 60$).

Figure~\ref{fig:ssl} shows the coding ceiling. Despite best-engineering SSL, coding at full labels ($\text{lf}=1.0$) reaches $\text{AUROC} \approx 0.45$--$0.47$---\emph{below} the simple \texttt{no\_svd\_lr} baseline (0.449). At low label fractions ($\text{lf}=0.05$), SSL v2 provides a modest gain (ssl\_v2\_r24 reaches 0.537 vs.\ no\_svd\_lr 0.488, $+4.9\text{pp}$), consistent with a regularization effect in the low-data regime. Science shows a larger SSL benefit ($+6.7\text{pp}$ at lf=0.05), consistent with the domain where features do carry signal.

\begin{figure*}[t]
  \centering
  \includegraphics[width=\textwidth]{fig_ssl_ceiling_v12.pdf}
  \caption{%
    \textbf{SSL ceiling in coding vs.\ science.}
    \textbf{(a)} In coding, all SSL variants and baselines cluster near AUROC $\approx 0.45$--$0.55$
    across all label fractions. Even the best-engineered SSL (ssl\_v2\_r24, VICReg, domain-specific
    encoder, 500 epochs) does not escape the ceiling at full labels.
    \textbf{(b)} In science, SSL v2 provides a meaningful $+6.7$pp gain at low label fractions,
    confirming that SSL can exploit latent structure when that structure exists.
    The contrast rules out label scarcity as the bottleneck in coding.
  }
  \label{fig:ssl}
\end{figure*}

\paragraph{Interpretation.} The SSL v2 result is the key experiment in this section. If there were latent structure in the coding feature space that correlates with correctness---even structure that linear probing cannot access with limited labels---SSL pre-training on 42,000 runs should find it. It does not. The training loss decreases (the encoder is learning \emph{something}), but what it learns is not predictive of correctness. The most parsimonious account: the token-trajectory features do not encode execution correctness, not because of data or model limitations, but because the two quantities are causally decoupled at the feature level.

\subsection{Method 5: Token-level surface-noise ablation (de-knotting)}
\label{sec:deknot}

A complementary hypothesis is that ``knot'' spans---circular reasoning segments where the model revisits the same idea without progress---add noise to the token stream, masking an otherwise-functional signal. If true, removing knot tokens should improve AUROC. We test this directly using a proper token-level ablation.

\paragraph{Protocol.} GLM-4-Flash (with domain-specific prompts) identifies knot character spans in $n=200$ balanced traces per domain (100 correct, 100 incorrect). We then map spans to token indices via the fast tokenizer's offset mapping and remove the corresponding positions from \emph{all five} per-token cached arrays (\texttt{tok\_conf}, \texttt{tok\_gini}, \texttt{tok\_logprob}, \texttt{tok\_neg\_entropy}, \texttt{tok\_selfcert}), as well as splicing the text. This is a full token-level ablation, not a text-only operation.

\paragraph{Results.} Figure~\ref{fig:deknot} shows knot prevalence and the AUROC response to removal.

\begin{figure*}[t]
  \centering
  \includegraphics[width=\textwidth]{fig_deknot_alldomains_v12.pdf}
  \caption{%
    \textbf{De-knotting ablation across all three domains.}
    \textbf{(a)} Knot prevalence differs by domain and by outcome.
    In math, knots are more prevalent in incorrect runs (82\% vs.\ 73\%)---the expected
    pattern if knots are failure indicators.
    In science, the difference is small and reversed (53\% incorrect vs.\ 60\% correct).
    In coding, knots are near-symmetric (90\% vs.\ 93\%).
    \textbf{(b)} Removing knot token spans from all five per-token signal arrays
    \emph{hurts} math AUROC ($\Delta=-0.049$) and leaves science and coding neutral.
    The math result confirms that knot tokens are the discriminative signal itself.
    The coding neutral result confirms that the failure is structural.
  }
  \label{fig:deknot}
\end{figure*}

\begin{itemize}
  \setlength\itemsep{0.15em}
  \item \textbf{Math}: 155/200 runs knotted; incorrect 82\% vs.\ correct 73\%. Removing all knot tokens \emph{hurts} AUROC from 0.502 to 0.453 ($\Delta = -0.049$). The knot-associated spans are the discriminative signal itself---they encode the failure pattern (persistent circular reasoning in incorrect runs). Removing them degrades the verifier.
  \item \textbf{Science}: 113/200 knotted; correct 60\% vs.\ incorrect 53\% (slight reversal). Removal neutral: AUROC $0.527 \to 0.521$ ($\Delta = -0.006$). Knot tokens carry near-zero signal.
  \item \textbf{Coding}: 183/200 knotted; near-symmetric (93\% vs.\ 90\%). Removal neutral: AUROC $0.460 \to 0.466$ ($\Delta = +0.006$). No masking strategy recovers signal, in any direction.
\end{itemize}

The math result is especially important: it rules out the interpretation that ``verifiers work in math because the knot spans there are informative.'' Rather, the knot spans in math \emph{are} the informative spans---removing them confirms it. The coding null result confirms that the feature-space failure is not surface noise.

\subsection{Convergence summary}

Table~\ref{tab:convergence} collects all five methods' results for coding.

\begin{table}[h]
\centering
\small
\begin{tabular}{lll}
\toprule
Method & Best coding AUROC & Verdict \\
\midrule
Interpretable probe (LogReg) & 0.434 (AoA) / 0.407 & Below baseline \\
Feature sweep (83 scalars) & 0.556 (top-4 combo) & Marginal \\
Nonlinear MLP & 0.499--0.507 & No change \\
SSL v1 (cross-domain)  & 0.502 ($r=16$, 100\% lf) & Collapse; no gain \\
SSL v2 (domain-specific, VICReg) & 0.537 (lf=0.05) / 0.454 (lf=1.0) & Low-label reg. only \\
De-knotting (token-level mask) & $\Delta = +0.006$ & Neutral \\
\bottomrule
\end{tabular}
\caption{All five methods applied to coding correctness prediction. No method escapes AUROC $\approx 0.5$ at full labels. SSL v2 shows a small regularization benefit at low label fractions but does not change the full-label ceiling.}
\label{tab:convergence}
\end{table}

% ─────────────────────────────────────────────────────────────────────────────
\section{Downstream Consequence: Best-of-$N$ Reranking}
\label{sec:reranking}
% ─────────────────────────────────────────────────────────────────────────────

The AUROC failure translates directly into a practical consequence. We run a best-of-$N$ selection experiment with $N=64$ candidates per problem: the run with the highest probe score at 100\% anchor is selected, and we report pass@1. The baseline is the per-problem mean accuracy (expected pass@1 under random selection).

\begin{figure}[h]
  \centering
  \includegraphics[width=\linewidth]{fig_reranking_v12.pdf}
  \caption{%
    Best-of-$N=64$ reranking. Probe reranking gains $+10.0\%$ on math and $+8.0\%$ on science,
    while yielding $-0.6\%$ (effectively zero) on coding. Error bars are 95\% bootstrap CIs
    from problem resampling ($B=10{,}000$, $n \ge 120$ problems per domain).
  }
  \label{fig:reranking}
\end{figure}

Table~\ref{tab:reranking} and Figure~\ref{fig:reranking} show the result. Math gains $+10.0$ percentage points, science gains $+8.0$ pp, and coding gains effectively nothing ($-0.6$ pp; 95\% CIs overlap completely). The measurement breakdown has a direct practical cost: deploying this probe for best-of-$N$ selection in coding wastes the selection budget and provides no lift.

\begin{table}[h]
\centering
\small
\begin{tabular}{lcccc}
\toprule
Domain & $n$ & Random & Probe & $\Delta$ \\
\midrule
Math    & 120 & .642 [.572, .712] & .742 [.658, .817] & \textbf{$+$.100} \\
Science & 198 & .602 [.550, .652] & .682 [.616, .748] & \textbf{$+$.080} \\
Coding  & 167 & .617 [.545, .689] & .611 [.539, .683] & $-$.006 \\
\bottomrule
\end{tabular}
\caption{Best-of-$N$ reranking pass@1 ($N=64$). Values in brackets are 95\% bootstrap CI ($B=10{,}000$). Coding pass@1 is statistically indistinguishable from random.}
\label{tab:reranking}
\end{table}

We do not claim that all verifiers fail on coding. Code-execution-aware verifiers, models with access to compiler feedback, or probes over hidden states have different information. The finding is specific to this feature family.

% ─────────────────────────────────────────────────────────────────────────────
\section{Why the Feature Family Fails in Coding}
\label{sec:why}
% ─────────────────────────────────────────────────────────────────────────────

\paragraph{The key asymmetry.} Math reasoning produces process-level evidence: when a solver is going in circles, the token stream says so---reflection counts increase, continuity drops, entropy oscillates. These patterns are causally downstream of the failure. In coding, the critical variable is \emph{whether the code compiles and passes tests}. A beautiful, well-organized CoT can produce incorrect code; a chaotic, uncertain CoT can produce correct code. The text is causally upstream of the code, but correctness is determined downstream of execution.

\paragraph{Why SSL does not help.} Self-supervised learning on CoT trajectories can learn structure that exists in those trajectories---temporal consistency, within-problem coherence, cross-anchor alignment. What it cannot do is discover a correlation between trajectory structure and correctness if that correlation does not exist in the data. The SSL result confirms this: loss goes down (structure is learned), but AUROC stays near 0.5 (correctness is not part of the learned structure).

\paragraph{Domain-conditioning as measurement invariance.} The psychometric concept of \emph{measurement invariance} \citep{vandenberg2000review} asks whether an instrument measures the same latent construct across groups. Our results suggest a measurement non-invariance: \texttt{traj\_reflection\_count} measures ``failure to converge'' in math (where convergence is defined by a closed-form answer) but measures ``exploratory effort'' in coding (where convergence requires successful execution). Same feature name, different construct.

\paragraph{What would fix coding.} Three directions are worth pursuing: (1)~\emph{Execution-aware features}: AST depth, cyclomatic complexity, or compiler-feedback signals that track executable state rather than text surface. (2)~\emph{Hidden-state probes}: intermediate activations may encode execution semantics not visible in the token stream. (3)~\emph{Multi-model replication}: separating model-specific effects from domain-specific effects requires testing the same features on a second model with different coding architecture.

% ─────────────────────────────────────────────────────────────────────────────
\section{Limitations and Conclusion}
\label{sec:limits}
% ─────────────────────────────────────────────────────────────────────────────

\paragraph{Limitations.} First, all experiments use a single model (DeepSeek-R1-0528-Qwen3-8B). The ceiling may differ for models with different reasoning styles. Second, the SSL experiments use up to 42,000 unlabeled runs---substantial but not exhaustive; at dramatically larger scale the picture may change. Third, the RL regime is excluded; verifier drift under RL fine-tuning is a different axis (regime-conditioning rather than domain-conditioning) and is studied separately.

\paragraph{Conclusion.} We have tested five methods on the question of whether token-trajectory features encode coding correctness in CoT traces. All five converge on no:

\begin{quote}
\textbf{Within this feature family, code correctness is not in the text.}
Simple probes, broad feature sweeps, nonlinear classifiers, self-supervised pre-training on 42,000 unlabeled runs, and token-level noise ablation all return the same result: coding AUROC cannot be meaningfully lifted above chance. The most informative negative result is the SSL experiment---if latent structure existed in these features, SSL would find it. It does not.
The contrast with math (AoA 0.958, $+10\%$ reranking gain) and science (AoA 0.799, $+8\%$ gain) makes the specificity of the coding failure clear. Before deploying a CoT-based verifier in a new domain, the right question is not ``what is the training accuracy?'' but ``does this feature family observe the right causal object?'' In coding, it does not.
\end{quote}

\paragraph{Artifacts.} Quantitative claims reproduce from \texttt{results/tables/}. Figure generation is in \texttt{scripts/gen\_figures\_v12.py}. Paper source is in \texttt{workshop/cotknot/paper/}.

% ─────────────────────────────────────────────────────────────────────────────
\bibliographystyle{plainnat}
\begin{thebibliography}{16}

\bibitem{wei2022cot}
Jason Wei et al.
\newblock Chain-of-thought prompting elicits reasoning in large language models.
\newblock In \emph{NeurIPS}, 2022.

\bibitem{wang2023selfconsistency}
Xuezhi Wang et al.
\newblock Self-consistency improves chain of thought reasoning in language models.
\newblock In \emph{ICLR}, 2023.

\bibitem{cobbe2021training}
Karl Cobbe et al.
\newblock Training verifiers to solve math word problems.
\newblock \emph{arXiv:2110.14168}, 2021.

\bibitem{lightman2023verify}
Hunter Lightman et al.
\newblock Let's verify step by step.
\newblock \emph{arXiv:2305.20050}, 2023.

\bibitem{wang2024mathshepherd}
Peiyi Wang et al.
\newblock Math-Shepherd: Verify and reinforce {LLM}s step-by-step without human annotations.
\newblock In \emph{ACL}, 2024.

\bibitem{zheng2025processbench}
Chujie Zheng et al.
\newblock ProcessBench: Identifying process errors in mathematical reasoning.
\newblock In \emph{ACL}, 2025.

\bibitem{zhang2025genrm}
Lunjun Zhang et al.
\newblock Generative verifiers: Reward modeling as next-token prediction.
\newblock In \emph{ICLR}, 2025.

\bibitem{lanham2023faithfulness}
Tamas K.\ Feher de Lanham et al.
\newblock Measuring faithfulness in chain-of-thought reasoning.
\newblock \emph{arXiv:2307.13702}, 2023.

\bibitem{liu2024intrinsic}
Dancheng Liu et al.
\newblock Large language models have intrinsic self-correction ability.
\newblock In \emph{NeurIPS}, 2024.

\bibitem{madaan2023selfrefine}
Aman Madaan et al.
\newblock Self-Refine: Iterative refinement with self-feedback.
\newblock \emph{arXiv:2303.17651}, 2023.

\bibitem{shinn2023reflexion}
Noah Shinn et al.
\newblock Reflexion: Language agents with verbal reinforcement learning.
\newblock \emph{arXiv:2303.11366}, 2023.

\bibitem{huang2024selfcorrect}
Jie Huang et al.
\newblock Large language models cannot self-correct reasoning yet.
\newblock \emph{arXiv:2310.01798}, 2023.

\bibitem{devlin2019bert}
Jacob Devlin et al.
\newblock {BERT}: Pre-training of deep bidirectional transformers for language understanding.
\newblock In \emph{NAACL}, 2019.

\bibitem{vandenberg2000review}
Robert E.\ Vandenberg and Charles E.\ Lance.
\newblock A review and synthesis of the measurement invariance literature.
\newblock \emph{Organizational Research Methods}, 3(1):4--70, 2000.

\bibitem{belinkov2022probing}
Yonatan Belinkov.
\newblock Probing classifiers: Promises, shortcomings, and advances.
\newblock \emph{Computational Linguistics}, 48(1):207--219, 2022.

\end{thebibliography}

\end{document}
